%====================================================================
% e3_identification_chain_level_platonic.tex
%
% Chain-Level E_3 Comparison with Explicit Associator Dependence
%
% Fixed-associator comparison of
% Z^{der}_{ch}(V_k(g)) and A^lambda.
% The strict object is a filtered chain model or a zig-zag over a
% cofibrant E_3 operad. A single strict map on the original BRST
% complex requires an additional rectification datum. Operad formality
% is used as an equivalence of deformation groupoids, not as an
% automatic algebra-level lifting theorem.
%
% Proved-here content:
% - thm:e3-identification-chain-level-associator-fixed
% Fixed-associator filtered chain model and strictness criterion.
% - prop:associator-dependence-explicit
% Different associators -> chain maps related by GRT_1 action;
% associated graded cohomology is unchanged.
% - thm:operad-level-to-algebra-level-lift
% Operad formality as MC-class comparison.
% - cor:e3-solvable-unconditional
% Solvable filtration criterion and abelian endpoint.
% - rem:grt1-action-on-identification
% GRT_1 acts on the space of chain-level identifications,
% trivially on associated graded cohomology, possibly nontrivially
% on cochains.
%
% Independent verification sources:
% Kontsevich (1999), Tamarkin (2003), Fresse Vol II (2017).
% Each source uses a distinct operadic model of E_n and a distinct
% formality route, so disjointness is structural, not merely formal.
%====================================================================

\providecommand{\Eone}{\mathsf{E}_1}
\providecommand{\Etwo}{\mathsf{E}_2}
\providecommand{\Ethree}{\mathsf{E}_3}
\providecommand{\En}{\mathsf{E}_n}
\providecommand{\Einf}{\mathsf{E}_{\infty}}
\providecommand{\Pone}{\mathsf{P}_1}
\providecommand{\Ptwo}{\mathsf{P}_2}
\providecommand{\Pthree}{\mathsf{P}_3}
\providecommand{\Pn}{\mathsf{P}_n}
\providecommand{\Assch}{\mathrm{Ass}^{\mathrm{ch}}}
\providecommand{\fg}{\mathfrak{g}}
\providecommand{\fh}{\mathfrak{h}}
\providecommand{\fz}{\mathfrak{z}}
\providecommand{\cA}{\mathcal{A}}
\providecommand{\cP}{\mathcal{P}}
\providecommand{\cQ}{\mathcal{Q}}
\providecommand{\GRT}{\mathrm{GRT}}
\providecommand{\GT}{\mathrm{GT}}
\providecommand{\MZV}{\mathrm{MZV}}
\providecommand{\Phi}{\Phi}
\providecommand{\KZ}{\mathrm{KZ}}
\providecommand{\Sym}{\operatorname{Sym}}
\providecommand{\HH}{\operatorname{HH}}
\providecommand{\CE}{\operatorname{CE}}
\providecommand{\Der}{\operatorname{Der}}
\providecommand{\Aut}{\operatorname{Aut}}
\providecommand{\End}{\operatorname{End}}
\providecommand{\Hom}{\operatorname{Hom}}
\providecommand{\id}{\mathrm{id}}
\providecommand{\MC}{\operatorname{MC}}
\providecommand{\Def}{\operatorname{Def}}
\providecommand{\coh}{\operatorname{coh}}

\chapter{Chain-Level $\Ethree$ Comparison with Explicit Associator Dependence}
\label{chap:e3-identification-chain-level-platonic}

\begin{remark}[$\Eone$-chiral input and $\Ethree$ output]
\label{rem:e3-id-platonic-ap-conformance}
\index{chiral notion discipline}
The chain-level $\Ethree$-identification of this chapter takes as
input the chiral algebra $V_k(\fg)$ in notion \textup{(A)} of
Warning~\ref{warn:multiple-e1-chiral} (a strict $\chirAss$-algebra
on the formal disk $D = \Spec\CC[[z]]$), with the homotopical
refinement to notion~\textup{(B)} ($A_\infty$-chiral structure in
$\End^{\mathrm{ch}}_{V_k(\fg)}$) implicit in the $A_\infty$
intertwiner of the formality lift
(Theorem~\ref{thm:operad-level-to-algebra-level-lift}). The
EK-quantum-vertex notion \textup{(C)} appears only through the
chiral quantum-group equivalence on the Koszul locus
in the ordered quantum-group chapter; notions \textup{(D)},
\textup{(E)} are not used.

The output is the $\Ethree$-structure on the derived chiral
center $Z^{\mathrm{der}}_{\mathrm{ch}}(V_k(\fg))$. Two
$\Ethree$-structures coexist:

\begin{enumerate}[label=\textup{(\Alph*)}]
\item \textbf{$\Ethree$-algebraic.} Higher Deligne center on the
 symmetric bar, requiring $\barB^{\Sigma}(A)$ as an
 $\Etwo$-coalgebra; no conformal vector; the lift route
 used in this chapter passes through an MC class in the
 deformation complex of a fixed cofibrant chain model of~$\Ethree$.
\item \textbf{$\Ethree$-topological.} Sugawara at non-critical
 level (proved in
 Chapter~\ref{chap:topologization-chain-level-platonic}); a
 conformal vector is required and the level must satisfy
 $k\ne -h^\vee$. Cohomological in general; chain-level on the
 original complex for class~$\mathbf{L}$
 (Theorem~\ref{thm:chain-level-E3-top-class-L}).
\end{enumerate}
Theorem~\ref{thm:e3-identification-chain-level-associator-fixed}
compares these structures after fixing an associator and a filtered
chain model. A strict map on the original BRST complex is asserted
only when the original-complex rectification datum is supplied; for
class~$\mathbf{L}$ this is exactly the condition verified by
Theorem~\ref{thm:chain-level-E3-top-class-L}.
\end{remark}

\begin{remark}[Families and obstruction spaces]
\label{rem:e3-id-family-enumeration}
\index{E3 identification!family enumeration}
The deformation space seen by $\Ethree$ formality is the space of
invariant symmetric bilinear forms on~$\fg$, not merely
$H^3(\fg;\CC)$. The Cartan map
$B\mapsto \omega_B(X,Y,Z)=B(X,[Y,Z])$ is an isomorphism for
semisimple~$\fg$ and has a kernel on the central or abelian part.
The comparison theorem therefore separates the families as follows.

\begin{center}
\renewcommand{\arraystretch}{1.2}
\small
\begin{tabular}{@{}lll@{}}
\toprule
Family & Chain-level status & Required datum \\
\midrule
$\fg$ simple
 & fixed-associator chain model
 & Cartan three-form and class-$\mathbf L$ rectification \\
$\mathfrak{gl}_N$
 & two-parameter comparison
 & $B_{\mathrm{tr}}$ and $B_{\mathrm{ab}}$ both matched \\
$\fg$ semisimple
 & componentwise comparison
 & one scalar per simple summand \\
$\fg$ reductive
 & conditional comparison
 & semisimple scalars plus central quadratic forms \\
$\fg$ abelian
 & explicit endpoint
 & invariant quadratic form only \\
$\fg$ nilpotent
 & filtered criterion
 & extension MC classes and invariant forms matched \\
$\fg$ solvable
 & filtered criterion
 & derived-series extensions split in the deformation groupoid
\end{tabular}
\end{center}

The word ``formality'' in the table means formality of the operad
together with equality of MC classes in the filtered deformation
groupoid. It does not mean that a quasi-isomorphism of associated
gradeds lifts automatically to a strict map of the original
algebras. The abelian endpoint is explicit. Nilpotent and solvable
algebras require the extension cocycles of the lower central or
derived filtration to be matched layer by layer
(Corollary~\ref{cor:e3-solvable-unconditional}).

For class~$\mathbf{M}$ (Virasoro, $\cW_N$): the chain-level
$\Ethree^{\mathrm{top}}$-structure on the original complex is
\textsc{Open}; weight-completed it is \textup{Proved}
(Volume~II, the $E_{k+2}^{\mathrm{top}}$ ladder theorem). For
class~$\mathbf{C}$ ($\beta\gamma$, $bc$), a chain-level statement
passes through the chosen bosonization model and inherits its
filtration.
\end{remark}

\begin{remark}[Bridge to the Climax Theorem anchors]
\label{rem:e3-id-bridge}
\index{Climax Theorem!E3 identification channel}
\index{universal holography!E3 identification channel}
The $\Ethree$-algebraic structure on
$Z^{\mathrm{der}}_{\mathrm{ch}}(V_k(\fg))$ is the higher Deligne
output of the Climax Theorem
(Chapter~\ref{chap:climax-platonic}): the bar differential as
$\mathrm{KZ}^*(\nabla^{\mathrm{Arnold}})$ supplies the
$\Etwo$-coalgebra structure on $\barB^{\Sigma}(V_k(\fg))$ from
which the higher Deligne center constructs the $\Ethree$-algebra
on $Z^{\mathrm{der}}_{\mathrm{ch}}(V_k(\fg))$. The chain map
$\varphi^{\Phi_{\mathrm{KZ}}}$ of
Theorem~\ref{thm:e3-identification-chain-level-associator-fixed}
is associator-rigidified. The localized bar--cobar and derived-center
statements are independent of this rigidification, while the chosen
chain representative is not. The
$\mathrm{GRT}_1$-torsor of
Remark~\ref{rem:grt1-action-on-identification} measures the gap.
The Vol~II universal-holography theorem
consumes the symmetric MC element on
$\Ethree$-rigidified $Z^{\mathrm{der}}_{\mathrm{ch}}$ to recover
the genus tower; conformance: this passage is via
Drinfeld center (right adjoint to forgetful), not via averaging.
\end{remark}

\section*{Fixed associator and obstruction tower}

Theorem~\ref{thm:e3-identification} of
Chapter~\ref{ch:en-koszul-duality} establishes the
$\Ethree$-deformation-theoretic identification
\[
 \{Z^{\mathrm{der}}_{\mathrm{ch}}(V_k(\fg))\}_{k + h^\vee}
 \;\simeq\;
 \{\cA^\lambda\}_{\lambda \in \lambda H^3(\fg)[[\lambda]]}
\]
as formal families, via $\En$-formality and an
obstruction-theoretic induction on~$\lambda = k + h^\vee$. The
argument identifies the two families as deformation classes, up
to $\Ethree$-quasi-isomorphism. What it does not exhibit is a
canonical strict chain map on the original dg complexes. At chain
level two extra choices enter: a cofibrant chain model for
$\Ethree$ and an associator rigidifying the formality
quasi-isomorphism.

This chapter proves the following precise replacement for the
naive strictness claim.

\begin{enumerate}[label=(\roman*)]
\item Fixing the Drinfeld KZ associator $\Phi=\Phi_{\KZ}$ gives a
filtered chain model and, under an explicit rectification
hypothesis, a strict filtered map
\[
 \varphi^{\Phi_{\KZ}} \colon
 Z^{\mathrm{der}}_{\mathrm{ch}}(V_k(\fg))
 \xrightarrow{\sim}
 \cA^\lambda
\]
of $\Ethree$-algebras compatible with the filtration by
$\lambda$-weight
(Theorem~\ref{thm:e3-identification-chain-level-associator-fixed}).
Without that rectification one obtains a filtered zig-zag through
cofibrant replacements.

\item We characterise the dependence of $\varphi^\Phi$ on the
chosen associator~$\Phi$:
two associators $\Phi, \Phi'$ produce chain maps
$\varphi^\Phi, \varphi^{\Phi'}$ related by the
Grothendieck--Teichm\"uller action on the chosen chain model.
For KZ and motivic associators the coefficients are periods in the
MZV algebra. On cohomology the action is unipotent and identity on
the associated graded; literal identity requires a specified
homotopy
(Proposition~\ref{prop:associator-dependence-explicit}).

\item We upgrade the operad-level formality input of
Proposition~\ref{prop:en-formality} to the correct algebra-level
criterion: formality identifies the filtered $\Ethree$ deformation
groupoid with the filtered $\Pthree$ deformation groupoid. A lift
exists exactly when the corresponding MC classes agree
(Theorem~\ref{thm:operad-level-to-algebra-level-lift}).

\item The finite derived series of a solvable algebra supplies a
finite filtration. It gives an unconditional comparison only in the
abelian endpoint and, more generally, when the extension MC classes
vanish or match in the filtered deformation groupoid
(Corollary~\ref{cor:e3-solvable-unconditional}).

\item We make the action of the prounipotent
Grothendieck--Teichm\"uller group $\GRT_1$ on the space of
chain-level rigidifications explicit: $\GRT_1$ acts on
$\cA^\lambda$ by filtered homotopy $\Ethree$-automorphisms. The
action on associated graded cohomology is trivial; faithfulness on
the concrete cochain model is a representation-theoretic condition,
not a formal consequence of Willwacher's theorem
(Remark~\ref{rem:grt1-action-on-identification}).
\end{enumerate}

The deformation class of Theorem~\ref{thm:e3-identification}
remains associator-independent. The chain representative is not.

\subsection*{Operadic formality inputs}
\label{sec:e3-identification-chain-level-hz4}

Three standard formality inputs are used, each with a different
chain model of little disks.

\begin{itemize}
\item \textbf{Path~$\alpha$:} Kontsevich~\cite{Kon99}. The
configuration-space operad of compactified Fulton--MacPherson
type, semi-algebraic chains, formality via explicit
integral-kernel formulae.

\item \textbf{Path~$\beta$:} Tamarkin~\cite{Tamarkin03}.
Graph-complex operadic model, $\Etwo$ formality via
Etingof--Kazhdan associator quantisation; the route emphasises
associator-dependence.

\item \textbf{Path~$\gamma$:} Fresse~\cite{Fresse17},
Vol.~II, Theorem~16.2.1 (and Theorem~16.1.1 for operadic
transfer). Chord-diagram and cooperadic models; intrinsic
formality of $\En$ for all $n \geq 2$ via mapping-space
homotopical obstruction theory.
\end{itemize}

A common source of dependence between these paths would
arise only if they collapsed to the same operadic model or to
the same formality cochain; neither collapse occurs.
Kontsevich uses semi-algebraic chains on
$\mathrm{FM}_n(\RR^n)$, Tamarkin uses Kontsevich's graph
complex with an associator input, and Fresse uses the
chord-diagram cooperad $\Lambda_n^*$. The three formality
quasi-isomorphisms agree on~$H^\bullet$ (forced by the
Cohen--Taylor computation of $H^\bullet(\En) = \Pn$) and
disagree on cochains by an element of $\GRT_1$, which is
precisely the content of the associator dependence analysed
below.

\section{Setup: derived chiral centre and \texorpdfstring{$\cA^\lambda$}{A lambda}}
\label{sec:e3-identification-setup}

Let $\fg$ be a finite-dimensional Lie algebra over~$\CC$ (or
over $\RR$ in the formality chain; we pass to~$\CC$ after the
formality input). Write~$V_k(\fg)$ for the universal affine
vertex algebra at level~$k$, on the formal disk
$D = \Spec\CC[[z]]$, with Sugawara critical level
$k_{\mathrm{crit}} = -h^\vee$ when $\fg$ is simple. Let
$\lambda = k + h^\vee$.

\subsection{The derived chiral centre as a filtered
\texorpdfstring{$\Ethree$}{E3}-algebra}

The derived chiral centre
\[
 Z^{\mathrm{der}}_{\mathrm{ch}}(V_k(\fg))
 \;:=\;
 C^\bullet_{\mathrm{ch}}(V_k(\fg), V_k(\fg))
\]
of chiral Hochschild cochains
(Construction~\ref{constr:chiral-hochschild-cochains}) carries
a natural filtered $\Ethree$-algebra structure after
topologization at non-critical level
(Theorem~\ref{thm:topologization}), and the chain-level refinement of
Chapter~\ref{chap:topologization-chain-level-platonic}
(Theorem~\ref{thm:topologization-chain-level-strict}) gives this
$\Ethree$-structure on the original BRST complex (not merely up
to quasi-isomorphism).

The filtration
\[
 F^p Z^{\mathrm{der}}_{\mathrm{ch}}(V_k(\fg))
 \;=\;
 \{\text{cochains of $\lambda$-weight } \geq p\}
\]
is complete and exhaustive. The associated graded is
\[
 \gr^\bullet Z^{\mathrm{der}}_{\mathrm{ch}}(V_k(\fg))
 \;\cong\;
 C^\bullet(\fg; \CC)
 \;=\;
 \Sym(\fg^\vee[-1])
 \;=:\;
 C^*(\fg),
\]
the Chevalley--Eilenberg complex of~$\fg$
(Theorem~\ref{thm:chiral-e3-cfg}(i)).

\subsection{The algebra \texorpdfstring{$\cA^\lambda$}{A lambda}
of shift-polynomials}

Dually, the Costello--Francis--Gwilliam
algebra~$\cA^\lambda$ (Theorem~\ref{thm:cfg}) is a filtered
$\Ethree$-algebra obtained by perturbative quantisation of
Chern--Simons theory on the formal disk. Its underlying
graded vector space is
\[
 \cA^\lambda
 \;=\;
 \Sym(\fg^\vee[-1])[[\lambda]]
 \;=\;
 C^*(\fg)[[\lambda]].
\]
Equivalently, $\cA^\lambda$ is the algebra of formal
power-series-valued shift-polynomials on~$\fg$: elements of
$C^*(\fg)$ whose coefficients are formal power series
in~$\lambda$. The $\Ethree$-structure deforms the
commutative-algebra structure on~$C^*(\fg)$ by a $\Pthree$
bracket at order~$\lambda$ proportional to the Killing form
$(-,-)$ on~$\fg$, followed by higher corrections.

\begin{notation}[Shift-polynomials]
\label{not:shift-polynomial}
We call~$\cA^\lambda$ the algebra of
\emph{$\lambda$-shift-polynomials} on~$\fg$: the underlying
module is $C^*(\fg)[[\lambda]]$, and the name records that the
$\Pthree$ bracket shifts by $\lambda$ at each order, with
coefficients controlled by invariant bilinear forms on~$\fg$.
\end{notation}

The filtration on $\cA^\lambda$ by $\lambda$-weight agrees
with the filtration on $Z^{\mathrm{der}}_{\mathrm{ch}}(V_k(\fg))$
at the level of associated gradeds, by
Theorem~\ref{thm:chiral-e3-cfg}(i).

\subsection{The problem}

Theorem~\ref{thm:e3-identification} produces an isomorphism of
\emph{deformation classes}: an equivalence
$Z^{\mathrm{der}}_{\mathrm{ch}}(V_k(\fg)) \simeq \cA^\lambda$
in the homotopy category of $\Ethree$-algebras over
$\lambda H^3(\fg)[[\lambda]]$. To obtain a chain map one must
rigidify the $\Ethree$-structure by choosing an associator; to
obtain a strict map on the original complex one must also rectify a
cofibrant chain model back to that complex. We make both requirements
explicit.

\section{The chain-level identification at fixed associator}
\label{sec:chain-level-at-fixed-associator}

\begin{theorem}[Fixed-associator $\Ethree$ comparison and strictness
criterion, $\Phi = \Phi_{\KZ}$]
\label{thm:e3-identification-chain-level-associator-fixed}
\ClaimStatusConditional
Let $\fg$ be a finite-dimensional Lie algebra over~$\CC$ and
let
\[
 \cZ_k(\fg):=Z^{\mathrm{der}}_{\mathrm{ch}}(V_k(\fg)),
 \qquad
 \cA^\lambda:=C^*(\fg)[[\lambda]]
\]
with their complete decreasing $\lambda$-weight filtrations.
Fix the Drinfeld Knizhnik--Zamolodchikov associator
$\Phi_{\KZ}$ and let
$C_*^{\Phi_{\KZ}}(\Ethree)$ denote the corresponding cofibrant
chain operad model of little $3$-disks. Suppose the following
three data are given.

\begin{enumerate}[label=\textup{(\roman*)}]
\item \textup{(Associated graded comparison.)}
There is an identification
\[
 \gr \cZ_k(\fg)\cong C^*(\fg)\cong \gr \cA^\lambda
\]
as $\Pthree$-algebras, and the induced $\Pthree$ brackets agree
on every component of
$\Sym^2(\fg^*)^\fg$. For simple~$\fg$ this is the single Cartan
three-form component; for reductive~$\fg$ it includes the central
quadratic forms.

\item \textup{(Filtered MC equality.)}
Under the formality equivalence of
Theorem~\ref{thm:operad-level-to-algebra-level-lift}, the two
filtered $\Ethree$-structures determine Maurer--Cartan elements
\[
 \theta_{\cZ},\theta_{\cA}
 \in
 \MC\bigl(F^1\Def_{\Pthree}(C^*(\fg))\bigr)
\]
which are gauge equivalent by an element of the pronilpotent group
\[
 \exp F^1\Def^0_{\Pthree}(C^*(\fg)).
\]

\item \textup{(Rectification datum.)}
A cofibrant replacement
$q\colon \cZ_k(\fg)^{\mathrm{cof}}\to \cZ_k(\fg)$ is fixed in
$C_*^{\Phi_{\KZ}}(\Ethree)$-algebras. If one wants a single strict
map on the original BRST complex, one must additionally give a
filtered $C_*^{\Phi_{\KZ}}(\Ethree)$-algebra rectification
identifying $\cZ_k(\fg)^{\mathrm{cof}}$ with $\cZ_k(\fg)$.
\end{enumerate}

Then there is a filtered zig-zag of
$C_*^{\Phi_{\KZ}}(\Ethree)$-algebras
\[
 \cZ_k(\fg)
 \xleftarrow{\;q\;}
 \cZ_k(\fg)^{\mathrm{cof}}
 \xrightarrow{\;\varphi^{\Phi_{\KZ}}\;}
 \cA^\lambda
\]
whose arrows are filtered quasi-isomorphisms and whose right arrow
reduces to the identity on $C^*(\fg)$ after passing to associated
graded. If the rectification datum in \textup{(iii)} is strict on
the original complex, the zig-zag contracts to a strict filtered
chain map
\[
 \varphi^{\Phi_{\KZ}}\colon
 \cZ_k(\fg)\longrightarrow \cA^\lambda
\]
of $C_*^{\Phi_{\KZ}}(\Ethree)$-algebras.

For simple~$\fg$ at non-critical level, conditions
\textup{(i)}--\textup{(ii)} are supplied by
Theorem~\ref{thm:e3-identification} and
Theorem~\ref{thm:chiral-e3-cfg}. The original-complex
rectification in \textup{(iii)} is supplied in class~$\mathbf L$
by Theorem~\ref{thm:chain-level-E3-top-class-L}; otherwise the
statement remains the zig-zag on a chosen chain model.

\noindent\textsc{Independent verification.}
\ClaimIndependentlyVerified[derived\_from=Decorator~$\gamma$
(Fresse~\cite{Fresse17}, Vol.~II, Theorem~16.2.1),
verified\_against=Decorator~$\alpha$ (Kontsevich
configuration-space formality~\cite{Kon99}),
disjoint\_rationale=Kontsevich semi-algebraic chains on
$\mathrm{FM}_n(\RR^n)$ vs.\ Fresse chord-diagram cooperad;
no common chain-level intermediate]
\end{theorem}

\begin{proof}
Fix the operad $C_*^{\Phi_{\KZ}}(\Ethree)$. The filtration is
complete, separated, and pronilpotent on the positive part
$F^1$ of the deformation complex, so the usual Deligne groupoid of
Maurer--Cartan elements is defined.

\smallskip
\noindent\textit{Step 1: formality as a deformation-groupoid
equivalence.}
Theorem~\ref{thm:operad-level-to-algebra-level-lift} identifies
filtered $C_*^{\Phi_{\KZ}}(\Ethree)$-deformations of
$C^*(\fg)$ with filtered $\Pthree$-deformations of $C^*(\fg)$,
on the level of MC groupoids. Thus the two chain models are
equivalent precisely when their $\Pthree$ MC elements are gauge
equivalent. This is assumption~\textup{(ii)}.

\smallskip
\noindent\textit{Step 2: exponentiating the gauge.}
Let $u\in F^1\Def^0_{\Pthree}(C^*(\fg))$ carry
$\theta_{\cZ}$ to~$\theta_{\cA}$. Since $F^1$ is pronilpotent,
$\exp(u)$ converges in the $\lambda$-adic topology and gives a
filtered isomorphism between the corresponding finite quotients
modulo~$F^{p+1}$ for every~$p$. Transporting this gauge through the
fixed formality equivalence gives a filtered
$C_*^{\Phi_{\KZ}}(\Ethree)$-morphism
\[
 \varphi^{\Phi_{\KZ}}_p\colon
 \cZ_k(\fg)^{\mathrm{cof}}/F^{p+1}
 \longrightarrow
 \cA^\lambda/F^{p+1}.
\]
The maps are compatible in~$p$ because $u$ lies in the complete
pronilpotent deformation complex.

\smallskip
\noindent\textit{Step 3: passage to the inverse limit.}
The finite quotient maps define
\[
 \varphi^{\Phi_{\KZ}}
 :=
 \varprojlim_p\varphi^{\Phi_{\KZ}}_p.
\]
The inverse limit is exact on this tower: each transition map is
surjective and each quotient $F^p/F^{p+1}$ is finite in every fixed
Chevalley--Eilenberg degree. This is the strong Mittag--Leffler
condition. The limit is a filtered
$C_*^{\Phi_{\KZ}}(\Ethree)$-morphism from the cofibrant model to
$\cA^\lambda$.

\smallskip
\noindent\textit{Step 4: quasi-isomorphism.}
On associated graded the map is the identity on~$C^*(\fg)$ by
assumption~\textup{(i)}. The spectral sequence of a complete
separated filtration therefore gives a quasi-isomorphism on the
underlying chain complexes. Compatibility with the
$C_*^{\Phi_{\KZ}}(\Ethree)$-operations is built into the
deformation-groupoid construction.

\smallskip
\noindent\textit{Step 5: strictness on the original complex.}
The left arrow $q$ is a weak equivalence by construction. A single
strict map out of $\cZ_k(\fg)$ is additional structure: it exists
when the chosen cofibrant model is rectified back to the original
BRST complex as a filtered
$C_*^{\Phi_{\KZ}}(\Ethree)$-algebra. This is exactly the
original-complex strictness condition separated in
Theorem~\ref{thm:topologization} and proved for class~$\mathbf L$
in Theorem~\ref{thm:chain-level-E3-top-class-L}. Without this
datum, the mathematically correct chain-level object is the
zig-zag.

\smallskip
\noindent\textit{Step 6: associator-dependence is captured in
Proposition~\ref{prop:associator-dependence-explicit}.}
The theorem fixes~$\Phi_{\KZ}$. Replacing it changes the
cofibrant operad model by the $\GRT_1$ action analysed below.

\smallskip
\noindent\textit{The binary test at $\mathfrak{sl}_2$.}
For $\fg = \mathfrak{sl}_2$, $h^\vee = 2$, and
Proposition~\ref{prop:e3-explicit-sl2} computes the leading
$\Pthree$ bracket
\[
 \{X,Y\}= (X,Y)/(k+2)
 \qquad X,Y\in \mathfrak{sl}_2.
\]
This binary operation is independent of the associator: associators
enter through arity at least~$3$. Hence the first test for the
comparison is the equality of this normalized Killing-form bracket;
the first possible associator-dependent correction lies in higher
arity and positive filtration.
\end{proof}

\section{Explicit associator dependence}
\label{sec:associator-dependence}

\begin{proposition}[Associator-dependence of the chain map]
\label{prop:associator-dependence-explicit}
\ClaimStatusConditional
Let $\Phi, \Phi' \in \GT_1(\CC)$ be two associators, and let
$\varphi^\Phi, \varphi^{\Phi'}$ be the chain maps of
Theorem~\ref{thm:e3-identification-chain-level-associator-fixed}
rigidified respectively by~$\Phi$ and~$\Phi'$ on chosen cofibrant
models. Then:

\begin{enumerate}[label=\textup{(\roman*)}]
\item \textup{(Quotient element.)}
The quotient $\sigma=\Phi'/\Phi \in \GRT_1(\CC)$ is a
prounipotent Grothendieck--Teichm\"uller element; its
logarithm
\[
 \log\sigma
 \;\in\;
 \mathrm{Lie}(\GRT_1)(\CC)
 \;=\;
 \widehat{\mathrm{grt}}_1
\]
lies in the completed Drinfeld--Ihara Lie algebra. For the KZ
associator and motivic associators its coefficients are periods in
the MZV algebra; an arbitrary complex associator has coefficients in
the chosen ground field.

\item \textup{(Action on $\cA^\lambda$.)}
Each element $\sigma \in \GRT_1(\CC)$ acts on $\cA^\lambda$
by a filtered homotopy
$C_*^\Phi(\Ethree)$-automorphism
\[
 \rho_{\cA}(\sigma)\colon \cA^\lambda \xrightarrow{\sim} \cA^\lambda
\]
after transport of operadic structure. The associated graded action
on the underlying $\Pthree$-algebra is the identity. In a graph model,
the positive-filtration terms are represented by Willwacher graph
cocycles \textup{(}Theorem~1.1 of~\cite{Willwacher15}\textup{)}.

\item \textup{(Explicit difference.)}
In the homotopy category of filtered chain
$\Ethree$-algebras,
\[
 \varphi^{\Phi'}
 \;\simeq\;
 \rho_{\cA}(\sigma) \circ \varphi^\Phi.
\]
After choosing the same cofibrant replacement and a representative of
$\rho_{\cA}(\sigma)$, this homotopy-class equality is represented by
a strict equality of chain maps.

\item \textup{(Cohomology.)}
The induced map
 $H^\bullet(\rho_{\cA}(\sigma)) \colon H^\bullet(\cA^\lambda)
\to H^\bullet(\cA^\lambda)$
preserves the induced filtration and is the identity on the
associated graded of cohomology. A literal identity on
$H^\bullet(\cA^\lambda)$ follows only after choosing a splitting or
a null-homotopy of the positive-filtration part. The deformation
class in the localized $\Ethree$-algebra category is independent of
the associator.

\item \textup{(Possible cochain-level detection.)}
If the Willwacher cocycle representing $\log\sigma$ acts nontrivially
on the chosen polydifferential model of~$\cA^\lambda$, then
$\varphi^{\Phi'}-\varphi^\Phi$ is nonzero on cochains. The first
possible associator-dependent term has arity at least~$3$ and
positive $\lambda$-filtration. For KZ-type choices the lowest odd
period that can occur is the $\zeta(3)$ wheel component
\textup{(}Furusho~\cite{Furusho11}\textup{)}.
\end{enumerate}

\noindent\textsc{Independent verification.}
\ClaimIndependentlyVerified[derived\_from=Brown~\cite{Brown12}
(motivic decomposition of $\mathrm{grt}_1$ via MZVs),
verified\_against=Furusho~\cite{Furusho11}
(pentagon equation rigidity for associators) and
Willwacher~\cite{Willwacher15}
(explicit graph-complex realisation of $\mathrm{grt}_1$ as
Kontsevich's polydifferential operators),
disjoint\_rationale=Brown works motivically in
category of mixed Tate motives over~$\Z$; Furusho works with
iterated integrals; Willwacher works with Kontsevich graphs.
The three sources agree on the abstract group but are
disjoint as proof routes]
\end{proposition}

\begin{proof}
Part (i) is the standard torsor property of associators: two
rigidifications differ by a prounipotent element of
$\GRT_1$. The MZV assertion is restricted to the KZ and motivic
subfamilies; Brown~\cite{Brown12} controls the motivic period
algebra, not the coefficients of an arbitrary complex associator.

Part (ii) is Willwacher's identification
$H^0(\mathrm{GC}_2)\cong\mathrm{grt}_1$
\cite{Willwacher15}. A graph cocycle changes the chain-level
formality morphism by a positive-filtration automorphism. Since
all such changes project to the same homology operad~$\Pthree$,
the associated graded $\Pthree$-algebra is unchanged.

Part (iii) is functoriality of the deformation-groupoid comparison
in the chosen associator. Transporting the operad model from
$\Phi$ to~$\Phi'$ is the same as post-composing by the automorphism
represented by~$\sigma=\Phi'/\Phi$. A homotopy-class equality is
all that is canonical. A strict equality requires the auxiliary
choice of the same cofibrant replacement and a chain representative
of~$\rho_{\cA}(\sigma)$.

Part (iv): a filtered automorphism that is the identity on associated
graded induces the identity on $\gr H^\bullet$, but it can still act
by an upper-triangular unipotent map on $H^\bullet$ before a splitting
is chosen. Thus the correct invariant statement is identity on
associated graded cohomology and equality in the localized
deformation class.

Part (v): associators do not affect binary operations. Their first
possible contribution appears in arity~$3$, through the graph-complex
classes representing $\mathrm{grt}_1$. Furusho's analysis of the
pentagon equation identifies the $\zeta(3)$ component for KZ-type
associators. Whether that class is detected on a particular
$\cA^\lambda$ cochain model is exactly the nonvanishing of the
corresponding Willwacher action.
\end{proof}

\begin{remark}[Scope of the nontriviality]
Part (v) of
Proposition~\ref{prop:associator-dependence-explicit} is
not a universal nonvanishing theorem. It says where a
chain-level difference can live. Detection on
$\cA^\lambda$ requires a nonzero representation of the relevant
graph cocycle on the chosen polydifferential model.
\end{remark}

\section{Operad formality and MC classes}
\label{sec:formality-lift}

The key technical tool, used in
Theorem~\ref{thm:e3-identification-chain-level-associator-fixed}
and in the analysis of
Proposition~\ref{prop:associator-dependence-explicit}, is the
formality comparison between filtered deformation groupoids.

\begin{theorem}[Operad formality as MC-class comparison]
\label{thm:operad-level-to-algebra-level-lift}
\ClaimStatusConditional
Let $n \geq 2$ and let $\cA, \cB$ be filtered
$\En$-algebras over~$\CC$ with complete exhaustive filtrations
such that $\gr^\bullet \cA \cong \gr^\bullet \cB$ as
$\Pn$-algebras. Fix an associator $\Phi \in \GT_1(\CC)$ and the
corresponding formality quasi-isomorphism
$C_*^\Phi(\En)\to \Pn$. Let
$\Def_{\Pn}(\gr\cA)$ denote the filtered deformation dg Lie algebra
of the common $\Pn$-algebra. Then:

\begin{enumerate}[label=\textup{(\roman*)}]
\item \textup{(Deformation groupoids.)}
The formality quasi-isomorphism induces an equivalence between
the filtered deformation groupoid of $C_*^\Phi(\En)$-algebra
structures on the common associated graded and the filtered
deformation groupoid
\[
 \MC\bigl(F^1\Def_{\Pn}(\gr\cA)\bigr)\big/\exp F^1\Def^0_{\Pn}(\gr\cA).
\]

\item \textup{(Lift criterion.)}
A quasi-isomorphism
$\psi\colon\gr^\bullet\cA\xrightarrow{\sim}\gr^\bullet\cB$
of $\Pn$-algebras lifts to a filtered
$C_*^\Phi(\En)$-quasi-isomorphism between cofibrant chain models
of~$\cA$ and~$\cB$ if and only if the two MC classes in the
groupoid of \textup{(i)} agree after transport by~$\psi$.
The lift is unique only up to gauge and higher homotopy.

\item \textup{(Original-complex strictness.)}
A single strict map
$\tilde\psi\colon \cA\to\cB$ on the original complexes is not a
formal consequence of operad formality. It requires a rectification
of the chosen cofibrant models back to the original filtered
algebras.

\item \textup{(Naturality in the associator.)}
Replacing~$\Phi$ by~$\Phi'$ transports the MC comparison by the
$\GRT_1$ element $\Phi'/\Phi$ and changes a chosen representative
by the Willwacher graph-complex action
\textup{(}Proposition~\ref{prop:associator-dependence-explicit}%
\textup{)}.

\item \textup{(Applicability to simple, reductive, and
solvable~$\fg$.)}
The associated-graded hypothesis is satisfied for
$\cA = Z^{\mathrm{der}}_{\mathrm{ch}}(V_k(\fg))$ and
$\cB = \cA^\lambda$ in the following scoped forms:
\begin{itemize}
\item $\fg$ simple: by
Theorem~\ref{thm:chiral-e3-cfg}(i), and the MC class is fixed by
the Cartan three-form;
\item $\fg$ reductive $\fg = \fg_{ss} \oplus \fz$: by the
K\"unneth identification
$\gr^\bullet \cA \cong C^*(\fg_{ss}) \otimes C^*(\fz)$
applied componentwise, provided all central quadratic-form
components are matched;
\item $\fg$ solvable: only under the finite-filtration criterion of
Corollary~\ref{cor:e3-solvable-unconditional}; the derived series
does not by itself kill extension MC classes.
\end{itemize}
\end{enumerate}

\noindent\textsc{Independent verification.}
\ClaimIndependentlyVerified[derived\_from=Fresse~\cite{Fresse17}
Vol.~II, Theorem~16.2.1 (mapping-space obstruction theory for
$\En$-algebras),
verified\_against=Tamarkin~\cite{Tamarkin03}
(Etingof--Kazhdan associator quantisation for $\Etwo$,
generalised to~$\En$ via Kontsevich--Soibelman),
disjoint\_rationale=Fresse works with chord-diagram cooperads
and completed mapping spaces;
Tamarkin works with graph-complex operads and EK quantisation;
the two routes share only the abstract statement of
operadic formality.
A third check: Kontsevich~\cite{Kon99}
semi-algebraic-chain formality at $n=2$ and its extension by
McClure--Smith~\cite{McClureSmith03} give the~$n=3$ case
independently]
\end{theorem}

\begin{proof}
Part (i): Fresse~\cite{Fresse17}, Vol.~II, Theorem~16.2.1,
applied to the operad quasi-isomorphism
$C_*^\Phi(\En)\to\Pn$, gives a Quillen equivalence on the
corresponding algebra categories. Passing to the formal
neighbourhood of the common associated graded gives the stated
equivalence of MC groupoids, as in
Lemma~\ref{lem:en-formality-deformation-classification}.

Part (ii): a filtered algebra structure is an MC element of the
positive-filtration deformation dg Lie algebra. A morphism lifting
$\psi$ exists exactly when the two MC elements become gauge
equivalent after transport by~$\psi$. If the MC classes differ, the
obstruction is precisely their difference in the Deligne groupoid;
operad formality cannot remove it. If they agree, choose a gauge
element in $\exp F^1\Def^0_{\Pn}(\gr\cA)$ and transport it through
the Quillen equivalence. The result is a quasi-isomorphism between
cofibrant models. Gauge choices give the higher homotopy
ambiguity.

Part (iii): the Quillen equivalence is a statement after replacing
objects by cofibrant models. A strict map on fixed original
complexes is a rectification problem. This is the strictness
distinction recorded in Theorem~\ref{thm:koszul-reflection}
\ref{KR-v} and in Theorem~\ref{thm:topologization}.

Part (iv): naturality is the functoriality of the Quillen
equivalence in the chosen formality morphism. The action of
$\Phi'/\Phi$ is the graph-complex action described in
Proposition~\ref{prop:associator-dependence-explicit}.

Part (v): for~$\fg$ simple, the associated graded is
$C^*(\fg)$ (Theorem~\ref{thm:chiral-e3-cfg}(i)), and the MC class
is determined by the Cartan three-form in
$H^3(\fg;\CC)$. For reductive
$\fg=\fg_{ss}\oplus\fz$, K\"unneth gives
$C^*(\fg)\cong C^*(\fg_{ss})\otimes C^*(\fz)$, while the central
summand contributes the kernel of the Cartan map
$\Sym^2(\fz^*)\to H^3(\fg)$. Hence all invariant quadratic-form
coordinates must be matched. For solvable~$\fg$, the finite
derived series gives a finite filtration but leaves extension MC
classes. Corollary~\ref{cor:e3-solvable-unconditional} gives the
exact criterion.

The comparison square used throughout is
\[
 \begin{tikzcd}
 \Hom_{\text{op}}(C_*(\En), \End_\cA)
 \ar[r, "\Phi_*"] \ar[d]
 & \Hom_{\text{op}}(\Pn, \End_\cA) \ar[d] \\
 \En\text{-}\Alg(\cA, \cB)
 \ar[r, "\tilde{\Phi}_*"]
 & \Pn\text{-}\Alg(\gr \cA, \gr \cB)
 \end{tikzcd}
\]
and it commutes up to the action of $\GRT_1$
(cf.\ Fresse~\cite{Fresse17}, Vol.~II, \S17.2).
\end{proof}

\section{Solvable filtrations}
\label{sec:solvable-unconditional}

\begin{corollary}[Solvable filtration criterion and abelian endpoint]
\label{cor:e3-solvable-unconditional}
\ClaimStatusConditional
Let $\fg$ be a finite-dimensional solvable Lie algebra
over~$\CC$. The derived series gives a finite filtration, but it
does not by itself produce a chain-level $\Ethree$ comparison.
The following criterion is necessary and sufficient for the
finite-filtration argument used in
Theorem~\ref{thm:e3-identification-chain-level-associator-fixed}.

\begin{enumerate}[label=\textup{(\roman*)}]
\item \textup{(Finite derived series.)}
The derived series
$\fg \supseteq \fg^{(1)} = [\fg,\fg]
\supseteq \fg^{(2)} = [\fg^{(1)},\fg^{(1)}]
\supseteq \cdots \supseteq \fg^{(d)} = 0$
is finite; each quotient $\fg^{(i)}/\fg^{(i+1)}$ is abelian.

\item \textup{(Tower of identifications.)}
The derived series induces a filtration
$F^\bullet \cA^\lambda$ and a compatible filtration on
$Z^{\mathrm{der}}_{\mathrm{ch}}(V_k(\fg))$ whose
associated graded at level~$i$ is
$C^*(\fg^{(i)}/\fg^{(i+1)}) = \Sym(\fg^{(i)}/\fg^{(i+1)})^\vee
[-1]$,
a graded-commutative polynomial/exterior algebra according to the
cohomological parity.

\item \textup{(Extension MC criterion.)}
Let
\[
 \epsilon_i(\cZ),\epsilon_i(\cA)
 \in
 H^1\Def_{\Pthree}\bigl(C^*(\fg^{(i)}/\fg^{(i+1)})\bigr)
\]
be the extension classes of the two filtered $\Pthree$-deformations
at the $i$th layer, including both the Lie-bracket extension cocycle
and the invariant bilinear-form component. The chain-level
comparison exists on the solvable filtration if and only if
\[
 \epsilon_i(\cZ)=\epsilon_i(\cA)
 \qquad\text{for all }i
\]
after transport through the associated-graded identifications.

\item \textup{(Strong convergence.)}
The filtration is finite
\textup{(}length~$d$\textup{)}, hence the tower of chain
identifications converges strongly whenever the extension criterion
in \textup{(iii)} holds. The limit
\[
 \varphi
 \;=\;
 \varphi_d \circ \varphi_{d-1} \circ \cdots \circ \varphi_1
\]
is a filtered chain-model quasi-isomorphism. Associator dependence
is absent only on layers whose $\Ethree$ structure is strictly
abelian.

\item \textup{(Abelian case.)}
For $\fg$ abelian (i.e.\ $d = 1$), $\cA^\lambda = C^*(\fg)
\otimes \CC[[\lambda]]$ with $\Pthree$ bracket determined solely
by the chosen invariant symmetric form on~$\fg$. If the same form
is used on both sides, the identification is the identity on the
underlying graded-commutative algebra.
\end{enumerate}

\noindent\textsc{Independent verification.}
\ClaimIndependentlyVerified[derived\_from=Direct inductive
argument on the derived series of~$\fg$ (elementary Lie
theory; Bourbaki~\cite{BourbakiLie},
Chap.~I, \S5, Theorem~1),
verified\_against=The filtered MC obstruction criterion of
Theorem~\ref{thm:operad-level-to-algebra-level-lift} and the
abelian Lie-cohomology computation (Weibel~\cite{WeibelHomological},
\S7.7),
disjoint\_rationale=Bourbaki supplies the finite filtration;
the MC criterion supplies the extension obstruction; Weibel supplies
the abelian endpoint]
\end{corollary}

\begin{proof}
Part (i) is classical: solvability is defined by finite
derived series (Bourbaki~\cite{BourbakiLie}, Chap.~I, \S5).

Part (ii): the PBW filtration on~$V_k(\fg)$ refines to a
filtration by the derived series of~$\fg$, because the
commutator relations in~$U(\fg)$ respect
$[\fg,\fg] \subset \fg$. Chiral Hochschild cochains inherit
this refinement, and the associated graded at level~$i$ is
the chiral Hochschild complex of the abelian quotient
$\fg^{(i)}/\fg^{(i+1)}$, which is
$C^*(\fg^{(i)}/\fg^{(i+1)})$ by
Theorem~\ref{thm:chiral-e3-cfg}(i) applied to the abelian
case.

Part (iii): a filtered extension of one abelian quotient by the
next is not determined by the associated graded alone. It is
determined by an MC class in the deformation complex of the graded
piece. This class records the Lie extension cocycle and the
invariant quadratic form that supplies the binary $\Pthree$
operation. Equality of these classes is exactly the lift criterion
of Theorem~\ref{thm:operad-level-to-algebra-level-lift} applied
layer by layer.

Part (iv): finite length removes analytic convergence issues. If the
extension MC classes match at every layer, each finite step supplies
a filtered quasi-isomorphism, and their finite composition gives the
claimed comparison. If one extension class differs, the obstruction
appears at the first layer where equality fails.

Part (v): for~$\fg$ abelian, the Lie bracket component of the MC
class is zero. The remaining $\Pthree$ bracket is the constant
biderivation attached to the chosen invariant symmetric bilinear
form. Matching that form makes the two MC classes equal and gives
the identity map on the underlying graded-commutative algebra.
\end{proof}

\begin{remark}[Heisenberg chiral algebra as a special case]
\label{rem:heisenberg-solvable-special-case}
The Heisenberg chiral algebra $H_k = V_k(\fh_1)$, where
$\fh_1$ is one-dimensional abelian, falls under the abelian
case of
Corollary~\ref{cor:e3-solvable-unconditional}(v): the
chain-level comparison is the identity once the level form is fixed,
and the $\Ethree$-structure is the constant biderivation determined
by that one-dimensional form.
This matches Theorem~\ref{thm:e3-identification}'s treatment
of class~G.
\end{remark}

\section{The GRT action on chain-level identifications}
\label{sec:grt1-action}

\begin{remark}[$\GRT_1$ action on the space of identifications]
\label{rem:grt1-action-on-identification}
The prounipotent Grothendieck--Teichm\"uller group~$\GRT_1(\CC)$
acts on the homotopy type of chain-level rigidifications in
Theorem~\ref{thm:e3-identification-chain-level-associator-fixed}
by post-composition:
\[
 \sigma \cdot \varphi
 \;=\;
 \rho_{\cA}(\sigma) \circ \varphi,
 \qquad
 \sigma \in \GRT_1(\CC),
\]
where $\rho_{\cA} \colon \GRT_1 \to \Aut^h(\cA^\lambda)$ is
Willwacher's realisation by graph-complex cocycles, as in
Proposition~\ref{prop:associator-dependence-explicit}(ii).

\begin{enumerate}[label=\textup{(\roman*)}]
\item \textup{(Torsor structure.)}
The set of chain-level $\Ethree$-quasi-isomorphisms
$\varphi^\Phi$ rigidified by an associator~$\Phi$ is a
$\GT_1(\CC)$-torsor after the cofibrant chain model and the
underlying deformation class have been fixed. Quotienting by this
torsor recovers the deformation-class identification of
Theorem~\ref{thm:e3-identification}.

\item \textup{(Associated graded cohomology.)}
For every $\sigma \in \GRT_1(\CC)$, the induced action on
$\gr H^\bullet(\cA^\lambda)$ is trivial
(Proposition~\ref{prop:associator-dependence-explicit}(iv)).
The unfiltered cohomology action is unipotent and becomes a literal
identity only after a splitting or null-homotopy is specified.

\item \textup{(Cochain detection.)}
Distinct associators can produce distinct cochain maps when the
corresponding Willwacher graph cocycle acts nontrivially on the
polydifferential model of~$\cA^\lambda$. This is a detection
statement about the representation of $\mathrm{GC}_2$, not a
formal faithfulness theorem for every~$\fg$.
For KZ-type associators, Brown's motivic MZV theorem
\cite{Brown12} explains the odd-zeta period coordinates that can
appear; it does not identify all of $\GRT_1(\CC)$ with the motivic
Galois group.

\item \textup{(Willwacher graph complex.)}
The Lie algebra $\widehat{\mathrm{grt}}_1$ is the degree-zero
cohomology of Kontsevich's graph complex~$\mathrm{GC}_2$
(Willwacher~\cite{Willwacher15}, Theorem~1.1), and each
generator acts on $\cA^\lambda$ by an explicit
polydifferential operator whose coefficients are integrals of
$\omega$-forms over configuration spaces on the upper
half-plane (Kontsevich graph integrals).

\item \textup{(No canonical choice.)}
There is no canonical choice of~$\Phi$ within~$\GT_1(\CC)$;
different canonical candidates include the KZ associator
$\Phi_{\KZ}$, the even-associator $\Phi_{\mathrm{ev}}$ of
Drinfeld, the Deligne associator, and the Alekseev--Torossian
Kashiwara--Vergne associator. Each gives a distinct
chain-level map
$\varphi^\Phi$, but the associated graded deformation class is the same.
\end{enumerate}

The role of this remark is to make
explicit the associator freedom that was left implicit in
Theorem~\ref{thm:e3-identification}. The deformation
identification is canonical after localization; the chain-level
rigidification is a torsor over~$\GT_1(\CC)$, with explicit
$\GRT_1$ action by
Willwacher graph cocycles.
\end{remark}

\section{The \texorpdfstring{$\mathfrak{sl}_2$}{sl2} binary bracket test}
\label{sec:russian-school-sl2}

We give an explicit computation at~$\fg = \mathfrak{sl}_2$
that verifies the associated-graded comparison and locates where
associator dependence can first occur.

\begin{proposition}[$\mathfrak{sl}_2$ associator-test]
\label{prop:sl2-associator-test}
\ClaimStatusConditional
For $\fg = \mathfrak{sl}_2$ with Chevalley basis $\{E, F, H\}$
and $h^\vee = 2$, the chain-level identification
$\varphi^{\Phi_{\KZ}}$ satisfies:

\begin{enumerate}[label=\textup{(\roman*)}]
\item \textup{(Order $\lambda^0$.)}
$\varphi^{\Phi_{\KZ}}_0$ is the identity on
$C^*(\mathfrak{sl}_2) = \Sym(\mathfrak{sl}_2^\vee[-1])
\cong \Lambda(e,f,h)$ with $|e|=|f|=|h|=1$.

\item \textup{(Binary $\Pthree$ bracket.)}
The binary operation on generators is
\[
\{X,Y\}
 =
 \frac{(X,Y)}{k+2},
 \qquad X,Y\in \mathfrak{sl}_2,
\]
with the normalized Killing form satisfying
$(E,F)=1$ and $(H,H)=2$. This is the same bracket computed in
Proposition~\ref{prop:e3-explicit-sl2}.

\item \textup{(Associator invisibility in arity two.)}
No associator changes the operation in \textup{(ii)}. The
associator acts through the arity-$3$ and higher part of the
operadic chain model. Thus the first possible difference between
$\varphi^{\Phi'}$ and $\varphi^{\Phi_{\KZ}}$ lies in positive
filtration and arity at least~$3$, represented in graph coordinates
by a Willwacher cocycle.

\item \textup{(Cohomological associated graded.)}
The induced action on
\[
 \gr H^\bullet(\mathfrak{sl}_2;\CC)
 =
 \CC\oplus \CC\cdot\omega_3,
 \qquad
 \omega_3=e\wedge f\wedge h/2,
\]
is the identity. Any unsplit cohomological action of an associator
is necessarily unipotent and invisible on this associated graded.
\end{enumerate}

\noindent\textsc{Independent verification.}
\ClaimIndependentlyVerified[derived\_from=Proposition~\ref{prop:e3-explicit-sl2}
and the normalized Killing-form computation,
verified\_against=CFG perturbative Chern--Simons binary bracket
\cite{CFG25} and the Wakimoto free-field realisation
\textup{(}Construction~\ref{constr:wakimoto-free-field}\textup{)},
disjoint\_rationale=The invariant-form computation is
Lie-theoretic; CFG uses perturbative Feynman graphs; Wakimoto uses
$\beta\gamma$ bosonisation]
\end{proposition}

\begin{proof}
Part (i) is immediate from
Theorem~\ref{thm:e3-identification-chain-level-associator-fixed}
on associated graded. Since the generators are odd, the symmetric
algebra on $\mathfrak{sl}_2^\vee[-1]$ is the exterior algebra
$\Lambda(e,f,h)$.

Part (ii) is Proposition~\ref{prop:e3-explicit-sl2}, equation
\eqref{eq:e3-p3-killing}, with $h^\vee=2$ and the trace-form
normalisation $(E,F)=1$, $(H,H)=2$. This computation is the
binary invariant-form component of the filtered MC class.

Part (iii): the associator is a solution of the pentagon equation
and controls parenthesisation in arity~$3$ and higher. The arity-$2$
homology of $\Ethree$ is already the binary piece of~$\Pthree$,
so no $\GRT_1$ element can alter the Killing-form bracket in
part~\textup{(ii)}.

Part (iv) follows from
Proposition~\ref{prop:associator-dependence-explicit}(iv):
the $\GRT_1$ action is positive-filtration and identity on
associated graded cohomology.
\end{proof}

\section{Fixed-associator consequences}
\label{sec:synthesis}

The results of this chapter assemble as follows:

\begin{enumerate}[label=\arabic*.]
\item \textbf{Fixed~$\Phi_{\KZ}$ comparison}
(Theorem~\ref{thm:e3-identification-chain-level-associator-fixed}):
the canonical chain-level object is a filtered zig-zag through a
cofibrant $\Ethree$ model. It becomes a strict map on the original
BRST complex only after the original-complex rectification datum is
supplied.

\item \textbf{Explicit associator-dependence}
(Proposition~\ref{prop:associator-dependence-explicit}): two
associators produce chain maps differing by a Willwacher
graph-complex $\GRT_1$-automorphism. KZ and motivic choices carry
MZV period coordinates. The associated graded cohomology map is the
identity; a literal cohomology identity requires a chosen splitting.

\item \textbf{Operad-level to algebra-level formality}
(Theorem~\ref{thm:operad-level-to-algebra-level-lift}): Fresse
Vol~II Theorem~16.2.1 gives an equivalence of deformation
groupoids. A lift exists exactly when the filtered MC classes agree.

\item \textbf{Solvable filtration criterion}
(Corollary~\ref{cor:e3-solvable-unconditional}): the finite derived
series removes inverse-limit convergence issues, but the extension
MC classes still have to match. The abelian endpoint is explicit.

\item \textbf{$\GRT_1$ torsor structure}
(Remark~\ref{rem:grt1-action-on-identification}): the space
of chain-level identifications is a $\GT_1(\CC)$-torsor,
trivial on associated graded cohomology and detectably nontrivial
on cochains only when the graph-complex action is nonzero.

\item \textbf{$\mathfrak{sl}_2$ binary test}
(Proposition~\ref{prop:sl2-associator-test}): explicit Wick
and invariant-form computation verify the binary $\Pthree$ bracket
at the smallest non-abelian example. Associators cannot change this
arity-$2$ operation.
\end{enumerate}

The deformation-theoretic identification
(Theorem~\ref{thm:e3-identification}) is thus sharpened into the
correct chain-level statement: deformation classes are
associator-independent after localization, while strict
representatives remember the associator and the rectification data.

\section{K3 Humbert comparison and the $\Ethree$ firewall}
\label{sec:kl-humbert-eighth-root}

The fixed-associator $\Ethree$ comparison is a local formal-disk
statement. It does not by itself prove a global Humbert monodromy
identity, a Lusztig root-of-unity equivalence, or a compact
Hall--Drinfeld recognition theorem. Those are arithmetic and
categorical comparison data living over the K3 period/Humbert
surface.

\begin{theorem}[Humbert--Lusztig--Mukai comparison firewall]
\label{thm:humbert-lusztig-mukai-8}
\ClaimStatusConditional
\index{Humbert divisor!monodromy order 8}
\index{Kazhdan--Lusztig!K3 BKM}
\index{Lusztig specialisation!$\ell = 8$}
Let the $\mathcal{B}$-family notation refer to the K3
Mukai/Borcherds comparison package of Volume~III. The
chain-level $\Ethree$ comparison of this chapter supplies none of
the following implications:
\[
 Z^{\mathrm{der}}_{\mathrm{ch}}\simeq_{\Ethree}\cA^\lambda
 \;\Longrightarrow\;
 \mathrm{ord}\bigl(\mathrm{Mon}_{H_1}(\mathcal L^{\Delta_5})\bigr)=8,
\]
\[
 Z^{\mathrm{der}}_{\mathrm{ch}}\simeq_{\Ethree}\cA^\lambda
 \;\Longrightarrow\;
 \ell_{\mathrm{Lusztig}}(\mathfrak u_{q_{\mathrm{K3}}})=8,
\]
or
\[
 Z^{\mathrm{der}}_{\mathrm{ch}}\simeq_{\Ethree}\cA^\lambda
 \;\Longrightarrow\;
 \mathrm{KL}_{q_{\mathrm{K3}}}\text{ is an equivalence.}
\]
The only scalar carried into this chapter from the K3 comparison
lane is the Mukai conductor normalisation
\[
 K^{\kappa_{\mathrm{ch}}}
 =
 2\,c_+(\mathrm{Mukai}(K3))
 =
 8,
\]
when that normalisation is supplied externally. The equality of this
integer with a Humbert monodromy order or a Lusztig order requires,
in addition, a holonomic $\mathcal D$-module
$\mathcal L^{\Delta_5}$ with normalized local monodromy, a
root-of-unity Hall--Drinfeld specialization, and a comparison map
identifying their Chern/root data. Under those extra data the scalar
check
\[
 \hbar^2 K^{\kappa_{\mathrm{ch}}}=-1
\]
is the identity $(-1/8)\cdot 8=-1$; it is not a proof of any
Kazhdan--Lusztig equivalence.
\end{theorem}

\begin{proof}[Proof]
The fixed-associator $\Ethree$ comparison is controlled by the
formal-disk deformation complex of $C^*(\fg)$ and by invariant
bilinear forms on~$\fg$. Its data are local in the chiral curve
direction. The monodromy of a holonomic $\mathcal D$-module on
$\overline{\mathcal{A}_2}$ is a global regular-singular
Riemann--Hilbert datum. The Lusztig order is a root-of-unity
specialisation datum in a quantum-group category. No functor in
Theorem~\ref{thm:e3-identification-chain-level-associator-fixed}
maps the first datum to either of the latter two.

The concordance makes the same separation: the proved Vol~I input at
the K3 base point is the Mukai-doubling conductor, while the Humbert
monodromy and Lusztig root-of-unity entries become comparable only
after the local monodromy normalisation, finite PBW truncation,
parity/root fixture, and comparison map to the K3
Hall--Drinfeld object are supplied
(Remark~\ref{rem:concord-2}). Thus the displayed implications are
invalid as consequences of this chapter. With all extra comparison
data installed, the remaining scalar identity is the arithmetic
evaluation $(-1/8)\cdot 8=-1$.
\end{proof}

\begin{remark}[Chain-level vs $(\infty,1)$-categorical status
at the Humbert wall]
\label{rem:humbert-chain-vs-infty}
Theorem~\ref{thm:humbert-lusztig-mukai-8} is stated and proved in
the Vol~I lane only as a firewall. A chain-level monodromy theorem
would require a Deligne canonical extension of
$\mathcal L^{\Delta_5}$ with residue computation at~$H_1$. An
$(\infty,1)$-categorical monodromy theorem would require the image of
a loop in
$\pi_1(\overline{\mathcal{A}_2}\setminus H_1)$ under the corresponding
Riemann--Hilbert functor. Both are legitimate mathematical lanes.
Neither is supplied by the local $\Ethree$ comparison.
\end{remark}

\begin{corollary}[Conditional Kazhdan--Lusztig target at
$q_{\mathrm{K3}} = e^{2\pi i/8}$]
\label{cor:kl-equivalence-k3-bkm}
\ClaimStatusConditional
\index{Kazhdan--Lusztig equivalence!K3 BKM}
\index{chiral quantum group!K3 Hall--Drinfeld double}
Assume that the K3 Hall--Drinfeld source has been recognized as the
compact object $\mathbf H_{\Delta_5}$, that the root-of-unity
specialisation at $\ell=8$ satisfies the PBW, radical-quotient, and
finite Frobenius-kernel hypotheses, and that the super
Etingof--Kazhdan/Kazhdan--Lusztig comparison functor is constructed
for the BKM superalgebra $\mathfrak g_{\Delta_5}$. Under these
hypotheses one obtains a candidate equivalence
of $(\infty,1)$-categories
\[
 \mathrm{Mod}^{\mathrm{ch}}(\mathbf{H}_{\Delta_5})^{\mathrm{loc,KL}}
 \;\simeq\;
 \mathrm{Mod}\bigl(\mathfrak{u}_{q_{\mathrm{K3}}}
 (\widehat{\mathfrak{m}}_{\mathrm{K3}})\bigr),
\]
with $q_{\mathrm{K3}} = e^{2\pi i / 8}$, where
$\mathfrak{u}_{q_{\mathrm{K3}}}$ is Lusztig's small quantum group for
the BKM superalgebra $\mathfrak{g}_{\Delta_5}$. This equivalence is
not a consequence of the local $\Ethree$ comparison in this chapter.
\end{corollary}

\begin{proof}[Proof sketch]
The hypotheses state exactly the data needed for such an equivalence:
a compact Hall--Drinfeld source, a finite root-of-unity quotient, and
a KL-type comparison functor for the relevant super BKM category. The
local $\Ethree$ theorem supplies none of these data. It is compatible
with them because its only scalar input on this lane is the external
normalisation $K^{\kappa_{\mathrm{ch}}}=8$.
\end{proof}

\begin{remark}[Why the conditional target uses $8$, not $12$]
\label{rem:why-eight-not-twelve}
A naive reader expects order $12$ (the order of a generic modular
monodromy around a cusp for an $\mathrm{SL}_2(\mathbb{Z})$ local
system; equivalently, the smallest $N$ with $\zeta_N \in K(\sqrt{-3})$).
The conditional K3 target uses $8$ only after the Humbert
$H_1$ normalisation and the Mukai conductor normalisation have been
identified by the external comparison map. The local $\Ethree$
comparison sees neither the elliptic cusp monodromy nor the Humbert
loop. It therefore cannot convert $12$ into $8$; it can only remain
compatible with the already supplied $8$-normalisation.
\end{remark}

\section*{Mathematical dependencies}
\label{sec:e3-chain-level-cross-refs}

This chapter is inserted into the characteristic-datum part
(see \ref{part:characteristic-datum} when defined) of Vol~I,
immediately after
Chapter~\ref{ch:en-koszul-duality}
and depends on the following mathematical inputs:

\begin{itemize}
\item The deformation-class identification
Theorem~\ref{thm:e3-identification}.
\item The two $\Ethree$ constructions and their associated gradeds:
Theorems~\ref{thm:e3-cs}, \ref{thm:cfg}, and
\ref{thm:chiral-e3-cfg}.
\item The topologization theorem and the class-$\mathbf L$
original-complex lift:
Theorems~\ref{thm:topologization} and
\ref{thm:chain-level-E3-top-class-L}.
\item The operadic formality input:
Proposition~\ref{prop:en-formality} and
Lemma~\ref{lem:en-formality-deformation-classification}.
\item The explicit binary test for $\mathfrak{sl}_2$:
Proposition~\ref{prop:e3-explicit-sl2}.
\end{itemize}

\endinput
